% Общая преамбула. Сборка: XeLaTeX (локально — tectonic).
% Оформление приближено к ГОСТ Р 7.0.11-2011: Times New Roman 14 пт,
% полуторный интервал, абзацный отступ 1,25 см, поля 30/15/20/20 мм.

% babel, а не polyglossia: русские подписи polyglossia задаёт через LICR
% (\CYRP\cyrr...), что под fontspec без T2A разваливается на \begin{document}.
\usepackage{fontspec}
\usepackage[english, russian]{babel}

\setmainfont{Times New Roman}[Ligatures=TeX]
\setmonofont{Menlo}[Scale=0.85]

\usepackage{geometry}
\geometry{a4paper, left=30mm, right=15mm, top=20mm, bottom=20mm}

\usepackage{setspace}
\onehalfspacing
\usepackage{indentfirst}
\setlength{\parindent}{1.25cm}

\usepackage{amsmath, amssymb}
\usepackage{graphicx}
\usepackage{booktabs}
\usepackage{longtable}
\usepackage{caption}
\captionsetup{labelsep=period, singlelinecheck=false, justification=centering}

\usepackage{listings}
\lstset{
  basicstyle=\ttfamily\small,
  breaklines=true,
  columns=fullflexible,
  keepspaces=true,
  showstringspaces=false,
  frame=single,
  framesep=4pt,
}

\usepackage[hidelinks]{hyperref}
\usepackage{cleveref}

% Нумерация формул, рисунков и таблиц — сквозная в пределах главы.
\numberwithin{equation}{chapter}
\numberwithin{figure}{chapter}
\numberwithin{table}{chapter}

% Заголовки глав по ГОСТ: прописными, по центру, без «Глава».
\usepackage{titlesec}
\titleformat{\chapter}[block]
  {\normalfont\bfseries\centering}{\thechapter}{1em}{\MakeUppercase}
\titlespacing*{\chapter}{0pt}{0pt}{24pt}

% Служебные макросы для терминологии, чтобы не расходиться по тексту.
\newcommand{\gist}{GiST}
\newcommand{\pg}{PostgreSQL}
\newcommand{\code}[1]{\texttt{#1}}
% Ссылка на коммит upstream: \commit{16fa9b2b30a}{Add support for building GiST index by sorting}
\newcommand{\commit}[2]{\code{#1} <<#2>>}
